
\documentclass{article}
\usepackage{colortbl}
\usepackage{makecell}
\usepackage{multirow}
\usepackage{supertabular}

\begin{document}

\newcounter{utterance}

\centering \large Interaction Transcript for game `multimodal\_referencegame', experiment `random\_grids', episode 0 with NEW\_Reflect\_Agent.
\vspace{24pt}

{ \footnotesize  \setcounter{utterance}{1}
\setlength{\tabcolsep}{0pt}
\begin{supertabular}{c@{$\;$}|p{.15\linewidth}@{}p{.15\linewidth}p{.15\linewidth}p{.15\linewidth}p{.15\linewidth}p{.15\linewidth}}
    \# & $\;$A & \multicolumn{4}{c}{Game Master} & $\;\:$B\\
    \hline

    \theutterance \stepcounter{utterance}  
    & & \multicolumn{4}{p{0.6\linewidth}}{
        \cellcolor[rgb]{0.9,0.9,0.9}{
            \makecell[{{p{\linewidth}}}]{
                \texttt{\tiny{[P1$\langle$GM]}}
                \texttt{You are given three images, one is called target and the other two are distractors.} \\
\texttt{Your task is to generate a referring expression that best describes the target image while distinguishing it from the two other distractor images.} \\
\texttt{The first image is the target, the second image is a distractor, and the third image is a distractor. Instruction: Describe the target image. Generate the referring expression starting with the tag "Expression: " for the given target image. Omit any other text.} \\
            }
        }
    }
    & \\ \\

    \theutterance \stepcounter{utterance}  
    & \multicolumn{4}{p{0.45\linewidth}}{
        \cellcolor[rgb]{0.95,0.95,0.95}{
            \makecell[{{p{\linewidth}}}]{
                \texttt{\tiny{[P1$|$internal]}}
                \texttt{Reflections from previous episodes:} \\
\texttt{{-}} \\
\texttt{ Failed. I did not generate a referring expression as instructed; I only acknowledged the task. Next time, I must directly output "Expression: " followed by the descriptive text, ignoring any preamble.} \\
\texttt{{-}} \\
\texttt{ Failed. I described the target's general scene but didn't distinguish it from the distractors, which also showed similar outdoor settings. Next time, I must identify and include a unique, distinguishing detail present only in the target (e.g., the specific activity or object arrangement).} \\
\texttt{{-}} \\
\texttt{ Failed. I did not generate a referring expression as instructed; I only acknowledged the task. Next time, I must directly output "Expression: " followed by a descriptive phrase that uniquely identifies the target image.} \\
            }
        }
    }
    & & \\ \\

    \theutterance \stepcounter{utterance}  
    & \multicolumn{4}{p{0.45\linewidth}}{
        \cellcolor[rgb]{0.95,0.95,0.95}{
            \makecell[{{p{\linewidth}}}]{
                \texttt{\tiny{[P1$|$internal]}}
                \texttt{{[}attempt 1 – valid{]}} \\
\texttt{Expression: The grid where the first row has 'â{-}č' in the first and last cells, and the second row has 'â{-}č' in the second, third, and fifth cells, while the third row has 'â{-}č' only in the second cell, and the fourth row has 'â{-}č' in the first, third, and fourth cells, and the fifth row has 'â{-}č' in the first, second, and third cells.} \\
            }
        }
    }
    & & \\ \\

    \theutterance \stepcounter{utterance}  
    & \multicolumn{4}{p{0.6\linewidth}}{
        \cellcolor[rgb]{0.9,0.9,0.9}{
            \makecell[{{p{\linewidth}}}]{
                \texttt{\tiny{[P1$\rangle$GM]}}
                \texttt{Expression: The grid where the first row has 'â{-}č' in the first and last cells, and the second row has 'â{-}č' in the second, third, and fifth cells, while the third row has 'â{-}č' only in the second cell, and the fourth row has 'â{-}č' in the first, third, and fourth cells, and the fifth row has 'â{-}č' in the first, second, and third cells.} \\
            }
        }
    }
    & & \\ \\

    \theutterance \stepcounter{utterance}  
    & & & \multicolumn{2}{p{0.3\linewidth}}{
        \cellcolor[rgb]{0.9,0.9,0.9}{
            \makecell[{{p{\linewidth}}}]{
                \texttt{\tiny{[GM$|$GM]}}
                \texttt{Expression: The grid where the first row has 'â{-}č' in the first and last cells, and the second row has 'â{-}č' in the second, third, and fifth cells, while the third row has 'â{-}č' only in the second cell, and the fourth row has 'â{-}č' in the first, third, and fourth cells, and the fifth row has 'â{-}č' in the first, second, and third cells.} \\
            }
        }
    }
    & & \\ \\

    \theutterance \stepcounter{utterance}  
    & & \multicolumn{4}{p{0.6\linewidth}}{
        \cellcolor[rgb]{0.9,0.9,0.9}{
            \makecell[{{p{\linewidth}}}]{
                \texttt{\tiny{[GM$\rangle$P2]}}
                \texttt{You are given three images.} \\
\texttt{You are also given a referring expression that describes one of the given images.} \\
\texttt{Your task is to select the image that matches the given referring expression.} \\
\texttt{Generate only the number (in text) of the image that matches the given expression by selecting first, second, or third.} \\
\\ 
\texttt{Expression: The grid where the first row has 'â{-}č' in the first and last cells, and the second row has 'â{-}č' in the second, third, and fifth cells, while the third row has 'â{-}č' only in the second cell, and the fourth row has 'â{-}č' in the first, third, and fourth cells, and the fifth row has 'â{-}č' in the first, second, and third cells.} \\
\texttt{Question: Which image does the expression refer to?} \\
\texttt{Start with the tag "Answer: ", followed by your selection. Omit any other text.} \\
            }
        }
    }
    & \\ \\

    \theutterance \stepcounter{utterance}  
    & & & \multicolumn{4}{p{0.45\linewidth}}{
        \cellcolor[rgb]{0.95,0.95,0.95}{
            \makecell[{{p{\linewidth}}}]{
                \texttt{\tiny{[P2$|$internal]}}
                \texttt{Reflections from previous episodes:} \\
\texttt{{-}} \\
\texttt{ Failed. The expression described a grid pattern, but I likely misinterpreted the special characters "â{-}ç" as literal rather than as placeholders for other symbols. Next time, I would interpret such characters as representing generic or unspecified content in the grid cells.} \\
\texttt{{-}} \\
\texttt{ Failed. I misinterpreted "first row, third column" as counting from the top{-}left, but the grid likely used a different coordinate system (e.g., bottom{-}left as origin). Next time, I must clarify the coordinate system (row/column indexing origin) before answering.} \\
\texttt{{-}} \\
\texttt{ Failed. I misinterpreted the expression, focusing on the presence of an "X" anywhere rather than its specific position in the *center cell*. I should have compared the exact grid coordinates of the "X" in each image to the description. Next time, I will parse spatial descriptions like "center cell" with strict, literal precision.} \\
            }
        }
    }
     \\ \\

    \theutterance \stepcounter{utterance}  
    & & & \multicolumn{4}{p{0.45\linewidth}}{
        \cellcolor[rgb]{0.95,0.95,0.95}{
            \makecell[{{p{\linewidth}}}]{
                \texttt{\tiny{[P2$|$internal]}}
                \texttt{{[}attempt 1 – valid{]}} \\
\texttt{Answer: first} \\
            }
        }
    }
     \\ \\

    \theutterance \stepcounter{utterance}  
    & & & \multicolumn{4}{p{0.6\linewidth}}{
        \cellcolor[rgb]{0.9,0.9,0.9}{
            \makecell[{{p{\linewidth}}}]{
                \texttt{\tiny{[GM$\langle$P2]}}
                \texttt{Answer: first} \\
            }
        }
    }
     \\ \\

    \theutterance \stepcounter{utterance}  
    & & & \multicolumn{2}{p{0.3\linewidth}}{
        \cellcolor[rgb]{0.9,0.9,0.9}{
            \makecell[{{p{\linewidth}}}]{
                \texttt{\tiny{[GM$|$GM]}}
                \texttt{first} \\
            }
        }
    }
    & & \\ \\

\end{supertabular}
}

\end{document}
